% Preface: the README's About the Course, Why native_sim, During the Course and Learning Outcomes
% sections, the course plan from info/README.md, and L01's prerequisites, as the front matter of a
% book.

\chapter{Preface}
\markboth{Preface}{Preface}

This book is for embedded C developers who want to build firmware with Zephyr. Every program in it
is built and run on Linux, on Zephyr's \code{native_sim} board, so no hardware is needed: not a
development board, not a debug probe, not a single wire.

\textbf{The hard part of Zephyr is not the kernel.} An engineer who has used FreeRTOS, or written a
scheduler, already knows what a thread, a mutex and a message queue are. What stops people is
everything around them: a workspace that west manages, a build that is CMake plus Kconfig plus
devicetree, and a driver model held together by macros. When Zephyr goes wrong, it is rarely
\code{k_thread_create()} at fault. It is \code{undefined reference to __device_dts_ord_42}, or a
build aborted by a Kconfig warning, or a misspelled \code{CONFIG_} macro that C quietly treats as
off.

So this book spends its first two chapters on the build and on devicetree, teaches the kernel as
\emph{the Zephyr way to do what you already know}, and finishes by writing a driver for a sensor
that does not exist, then testing and debugging it. Along the way it covers:
\begin{itemize}
  \item The west workspace, the manifest, and pinning a Zephyr release.
  \item The build: CMake, \code{prj.conf}, Kconfig, and what to do when the build rejects an
    option.
  \item Logging.
  \item Devicetree: nodes, properties, bindings, overlays, and the macros that read them at
    compile time.
  \item Getting devices from the devicetree, and GPIO through Zephyr's emulated GPIO controller.
  \item Threads, priorities, cooperative and preemptive scheduling, and time.
  \item Interrupt context, and deferring work out of it with work queues.
  \item Semaphores, mutexes with priority inheritance, message queues, events and polling.
  \item Application architecture with zbus and the State Machine Framework.
  \item The device driver model: API structs, instances, init levels and bindings.
  \item Emulators, which stand in for hardware in drivers and tests.
  \item Testing with ztest and Twister.
  \item The shell, settings stored on a simulated flash, and networking through host sockets.
  \item Debugging on the host: GDB, sanitizers, fatal error handling and coverage.
\end{itemize}

\section*{Why native\_sim}
\code{native_sim} compiles Zephyr, kernel and application together, into an ordinary Linux
executable. There is no instruction set simulator: your code runs natively on the host CPU, and the
hardware around it, a timer, an interrupt controller, a GPIO controller, an I2C bus, a flash chip,
is simulated. The book uses it for three reasons.

\textbf{Everything builds and runs on a laptop.} No boards to hand out, no probes, no drivers for a
USB debugger in WSL. A build takes seconds and a run takes less.

\textbf{It is deterministic.} Time on \code{native_sim} is simulated. A program that prints the
uptime prints the same numbers on every run and on every machine, which is what lets the exercises
compare your output with the expected output line by line.

\textbf{The host's tools work on it.} GDB, AddressSanitizer, UndefinedBehaviorSanitizer, gcov and
Valgrind all work on a \code{native_sim} program as on any other Linux program.

The counterweight belongs on the first page: \textbf{\code{native_sim} is not a microcontroller.}
Code runs in zero simulated time, so an interrupt cannot preempt a running thread until that thread
calls into the kernel. A busy loop waiting for an interrupt never ends. Nothing here teaches pin
multiplexing, clocks, power or flashing. Chapter~\ref{c3:ch} shows the time model on purpose,
because knowing what the simulator does \emph{not} do is as important as knowing what it does.

\section*{What you will do}
Along the way you will write code, not just read it. You will:
\begin{itemize}
  \item Set up a pinned Zephyr workspace and build and run applications on \code{native_sim}.
  \item Configure applications with Kconfig and describe hardware with devicetree overlays.
  \item Write multi-threaded applications with the kernel's scheduling and IPC primitives.
  \item Structure an application around zbus channels and a state machine.
  \item Write a devicetree-driven sensor driver, and an emulator to test it against.
  \item Test with ztest and Twister, and debug with GDB and the sanitizers.
\end{itemize}

The examples and exercises stay close to typical embedded problems: blinking an LED and reading a
button through an emulated GPIO controller, periodic sampling with timers and work queues, producer
and consumer threads, priority inversion produced and then fixed, a driver for an I2C temperature
sensor, and a shell command and a setting that survives a restart.

\section*{What you should be able to do at the end}
Having worked through the book, you should be able to:
\begin{itemize}
  \item Set up a Zephyr workspace and explain what west, CMake, Kconfig and devicetree each
    contribute.
  \item Find out why a configuration option or devicetree node did not take effect.
  \item Read a devicetree and binding, and access devices and properties from C.
  \item Choose between cooperative and preemptive threads, and between an ISR, a work item and a
    thread.
  \item Use Zephyr's IPC primitives, and explain when each one is the right tool.
  \item Write a devicetree-instantiated driver for Zephyr's device model.
  \item Test drivers and applications with ztest, emulators and Twister.
  \item Explain what \code{native_sim} simulates faithfully and what it does not.
\end{itemize}

\section*{What you are assumed to know}
The subject of this book is \textbf{Zephyr}, not embedded programming in general. You are expected
to arrive already comfortable with:
\begin{itemize}
  \item C programming: functions, pointers, structs and the rest of the language.
  \item A Linux terminal, \code{make} and, at the level of reading it, CMake.
\end{itemize}
Basic RTOS concepts, threads, priorities and blocking, help, but Chapters~\ref{c3:ch}
and~\ref{c4:ch} cover them again in Zephyr's terms. What is underneath them, the context switch
itself, is the subject of a different course, \emph{Real-Time Kernel Design}.

\section*{How the book is organised}
The book is typeset from a six-lecture course, and each chapter is one lecture. Its sections are
the lecture's appendices, in the same order and with the same subsections, so a chapter here and a
lecture directory in the repository can be read side by side:

\begin{booktable}{@{}p{5.2em}p{3.6em}L@{}}
  \thead{Chapter} & \thead{Lecture} & \thead{Topic} \\
  \midrule
  1 & \file{L01} & Zephyr, west and the build \\
  2 & \file{L02} & Devicetree \\
  3 & \file{L03} & Kernel I: threads, time and work queues \\
  4 & \file{L04} & Kernel II: communication and application architecture \\
  5 & \file{L05} & Writing a driver \\
  6 & \file{L06} & Subsystems and debugging \\
\end{booktable}

Every chapter opens with what is in it, works through the material, closes with a review of what
you should now be able to explain, and ends with its exercises.

\textbf{The live-coding examples are not printed.} Each chapter is taught from a handful of small
applications, which live in the companion repository under \file{lectures/LNN/examples}; every one
has a README saying how to build and run it and what to look at. The book names them where they
belong, and the code in the text is the code that makes a point, not a listing of every file.

\section*{What the exercises look like}
The exercises come in \textbf{sets}. A set builds one small application, or one part of a larger
one, and each exercise in a set adds a piece of it. Every exercise is labelled with its kind:

\begin{booktable}{@{}p{7.4em}L@{}}
  \thead{Kind} & \thead{What it asks of you} \\
  \midrule
  \kindtag{Code} & Write C, build it for \code{native_sim}, and check it with \code{make test},
    which also compares what it printed with the output the exercise gives. \\
  \kindtag[fillaccent2]{Reflection} & Answer in prose: explain what your program does, or why the
    build behaved as it did. \\
\end{booktable}

\textbf{The solutions are not in this book.} They are published in the companion repository, under
\file{lectures/LNN/appendix/solutions}, a few days after each lecture. Try each exercise before you
read them.

\section*{Two written papers, and what they are not}
Nothing in the course this book comes from is marked. Assessment is the exercise set after every
chapter, each with a published solution, and the questions in every chapter's review.

Appendices~\ref{exam:ch:about} to \ref{ansb:ch:answers} hold two three-hour papers with worked
solutions, and they check something else: \textbf{what you can reconstruct on paper, with nothing
in front of you.} Eight questions each, mixing what a build decided and why with the artifacts a
Zephyr application is made of, and no build in the room to tell you whether you are right.

\textbf{They exist purely so that you can test your own knowledge after working through the book.
They gate nothing, they are not a qualification, and no part of the book requires them.}

\textbf{Take one after the last chapter.} Both papers draw on all six chapters, so sitting one
partway through examines material you have not read yet.

\section*{Setting up}
Everything in this book runs in a Linux terminal: on Linux itself, or on Windows through WSL 2.
\Secref{c1:app:a} installs the packages, sets up the workspace and checks that it works; do that
first, because fetching Zephyr takes a few minutes. In short, it is:

\begin{shell}
git clone https://github.com/qrtech-academy/zephyr.git
cd zephyr
make setup
\end{shell}

\noindent\code{make setup} creates a Python environment in \file{.venv}, installs west into it, and
fetches the Zephyr release the manifest pins into \file{deps/}. Nothing is installed outside the
repository. Then, in every new terminal:

\begin{shell}
source env.sh
\end{shell}

\noindent Three more commands cover everything the book asks of the repository:

\begin{shell}
# Build an example, and run it.
west build -p -b native_sim/native/64 lectures/L01/examples/hello
west build -t run

# Test your exercises.
make test
\end{shell}

\noindent You write your exercise solutions in the repository, in \file{lectures/LNN/exercises},
where each lecture's README says which file goes where. \code{make test} then runs every lecture's
tests against them: a set you have not started is reported as \code{SKIP}, with the file its tests
are waiting for, and a set you have written is built, run and compared with the output the exercise
gives.
